\section{Conserved Payloads and Exit Channels}
\label{sec:decay_concepts}
\subsection{Fragmentation and conservation}
A parent and its complete exit cohort are two representations of a conserved payload. For muon decay, the electron and neutrino products together conserve charge and four-momentum. Angular-momentum conservation includes orbital contributions; spin labels are not ordinary additive scalar checksums. Calling neutrinos lattice phonons is a proposed interpretation, not a consequence of energy conservation or a Landauer bound.

\subsection{Partial widths and branching fractions}
For exclusive channels with $\Gamma_i\geq0$, adopt
\begin{equation}
\Gamma_{\rm tot}=\sum_i\Gamma_i,\qquad
 B_i=\frac{\Gamma_i}{\Gamma_{\rm tot}}.
\end{equation}
If $\Gamma_{\rm tot}>0$, the branching fractions sum to one. In an independent-hazard model, survival factors multiply and hazards add. Lean checks these implications. A multiplicity $C$ multiplies one width by $C$ only when those channels have equal partial widths. Counting labels does not prove that equality or determine hadronic matrix elements.

\subsection{Reachability and waiting times}
An allowed route and a nonzero rate are separate inputs. The finite descent argument bounds the number of transitions once they occur; it does not determine the waiting times between them. A deterministic update may remain in a state indefinitely unless further dynamics causes an exit.
